\documentclass[11pt,a4paper]{ctexart}

% Shared preamble for the Lena Voita seq2seq-and-attention Chinese lesson.
\usepackage[margin=2.2cm]{geometry}
\usepackage{amsmath,amssymb,amsthm,mathtools}
\usepackage[dvipsnames]{xcolor}
\usepackage{hyperref}
\usepackage{enumitem}
\usepackage[most]{tcolorbox}
\usepackage{tikz}
\usepackage{listings}
\usepackage{booktabs}
\usepackage{array}
\usepackage{caption}
\usetikzlibrary{arrows.meta,calc,positioning,fit,matrix}

\hypersetup{colorlinks=true,linkcolor=blue!55!black,urlcolor=blue!55!black}
\setlist[enumerate]{itemsep=0.35em, topsep=0.45em}
\setlist[itemize]{itemsep=0.35em, topsep=0.45em}

\newtheorem{theorem}{命题}[section]
\newtheorem{definition}[theorem]{定义}
\newtheorem{example}[theorem]{例}
\newtheorem{remark}[theorem]{备注}

\tcbset{
  commonbox/.style={
    enhanced, breakable,
    boxrule=0.8pt, arc=2mm,
    left=1.4mm, right=1.4mm, top=1mm, bottom=1mm,
    colback=white
  }
}
\newtcolorbox{problemblock}{commonbox,colframe=BrickRed!65!black,colback=BrickRed!4,title=背景问题,fonttitle=\bfseries}
\newtcolorbox{intuitionblock}{commonbox,colframe=RoyalBlue!70!black,colback=RoyalBlue!4,title=朴素直觉,fonttitle=\bfseries}
\newtcolorbox{mathbox}[1]{commonbox,colframe=Purple!60!black,colback=Purple!4,title=#1,fonttitle=\bfseries}
\newtcolorbox{engbox}{commonbox,colframe=ForestGreen!60!black,colback=ForestGreen!4,title=代码 / 工程对应,fonttitle=\bfseries}
\newtcolorbox{keypointblock}{commonbox,colframe=black!65,colback=black!3,title=记住这一点,fonttitle=\bfseries}

\lstset{
  basicstyle=\ttfamily\small,
  breaklines=true,
  columns=fullflexible,
  keywordstyle=\color{Purple!70!black}\bfseries,
  commentstyle=\color{gray!70!black}\itshape,
  stringstyle=\color{ForestGreen!50!black},
  showstringspaces=false,
  language=Python,
  frame=single,
  rulecolor=\color{black!20},
  framesep=2pt,
  xleftmargin=0.5em
}

\newcommand{\vect}[1]{\boldsymbol{#1}}
\newcommand{\E}{\mathbb{E}}
\newcommand{\R}{\mathbb{R}}
\newcommand{\softmax}{\mathrm{softmax}}
\newcommand{\score}{\mathrm{score}}


\title{Seq2seq 与 Attention\\\large 对照 Lena Voita NLP 课程的零基础精读}
\author{}
\date{\today}

\begin{document}
\maketitle

\begin{abstract}
本文是对 Lena Voita 的 NLP 课程章节
\textit{Sequence to Sequence (seq2seq) and Attention} 的零基础中文精读。
按原文目录顺序，把每一节都拆开重讲：
先抛出\emph{背景问题}，再给\emph{朴素直觉}，
然后写成\emph{数学表达}，最后给出\emph{代码 / 工程对应}。
读者只需要会基本的线性代数和 softmax 即可入门。
\end{abstract}

\tableofcontents
\newpage

\section{序列到序列基础 (Sequence to Sequence Basics)}
\label{ch:seq2seq-basics}

\subsection{章节定位与最该问的几个问题}

很多任务的输入和输出都是\emph{长度可变}的序列：把一句英文翻译成中文、把一段长文压缩成摘要、把一段
Python 代码翻译成等价的 C++、甚至把一张图片描述成一句话——这些都属于
\emph{序列到序列} (sequence to sequence, seq2seq) 任务。它们共享同一个抽象：给定源序列
$\vect{x} = (x_1, \dots, x_S)$，要生成目标序列 $\vect{y} = (y_1, \dots, y_T)$，而 $S$ 和 $T$
通常\emph{不相等}、也不预先固定。

\begin{problemblock}
在正式进入公式之前，先把本章想回答的几个朴素问题列出来。读完这一章你应该可以回答：
\begin{enumerate}
    \item Encoder 和 decoder 分别是干什么的？为什么要拆成两块？
    \item “条件语言模型” (conditional language model, CLM) 跟普通的语言模型差在哪？
    \item 训练时为什么用\emph{交叉熵}损失？它和最大似然是同一件事吗？
    \item 训练时 decoder 看到的“上一个词”到底是真值还是模型自己的预测？为什么这么做？
    \item 推理 (inference) 时为什么 \emph{greedy} 不够，要用 \emph{beam search}？
    \item Beam 越大就越好吗？
\end{enumerate}
\end{problemblock}

本章只覆盖最朴素的 seq2seq 框架（两个 RNN 之间用一个固定向量传话）。这个框架有一个明显的短板：
信息全靠一个固定向量\emph{瓶颈}，长句子会“记不住”。第二章引入的 \textbf{attention} 正是为了
解决这个瓶颈。所以请把这一章当成“为什么我们后面非得有 attention 不可”的铺垫。

\subsection{Encoder–Decoder 框架}
\label{sec:encdec}

\begin{intuitionblock}
最朴素的直觉：既然输入输出都是序列，而且通常\emph{长度不同}，那就分两步走。
\begin{itemize}
    \item \textbf{Encoder}（编码器）：把源序列 $\vect{x}$ \emph{读完}，压缩成某种“它到底在说什么”
          的内部表示。
    \item \textbf{Decoder}（解码器）：拿到这个表示，\emph{一个词一个词地写出} 目标序列 $\vect{y}$。
\end{itemize}
\end{intuitionblock}

为什么这种“先读后写”的拆分如此自然？
\begin{itemize}
    \item 输入和输出长度通常不同，无法用单一前馈网络一一对应。
    \item 在写第一个目标词之前，最好已经看完整个源句——否则连主语是谁都还不知道。
    \item 拆开之后，encoder 和 decoder 可以使用\emph{不同}的网络结构（RNN、CNN、Transformer
          都可以），便于实验和组合。
\end{itemize}

\begin{remark}
此后所有更复杂的模型——加了 attention 的 RNN、Transformer、甚至现代 LLM 中的
encoder-only / decoder-only 变体——都没有抛弃这个骨架，而只是在改“encoder 怎么编”、
“decoder 怎么用编码出来的东西”这两件事。
\end{remark}

\medskip
下面这张极简的“三盒图”足够你在脑子里画：

\begin{center}
\begin{tikzpicture}[
    box/.style={draw, rounded corners, minimum width=2.4cm, minimum height=1cm, align=center},
    arr/.style={-{Latex}, thick}
]
    \node[box, fill=blue!8]   (enc) at (0,0)   {Encoder \\ (读 $\vect{x}$)};
    \node[box, fill=purple!8] (rep) at (4.2,0) {内部表示 \\ $\vect{h}$};
    \node[box, fill=green!8]  (dec) at (8.4,0) {Decoder \\ (写 $\vect{y}$)};
    \draw[arr] (enc) -- (rep);
    \draw[arr] (rep) -- (dec);
    \node[below=0.05cm of enc] {\small 源序列 $\vect{x}$};
    \node[below=0.05cm of dec] {\small 目标序列 $\vect{y}$};
\end{tikzpicture}
\end{center}

\subsection{条件语言模型 (Conditional Language Model)}
\label{sec:clm}

\begin{intuitionblock}
普通的\emph{语言模型} (LM) 干的事是：在没有任何外部输入的情况下，告诉你下一个词最可能是什么。
\emph{条件}语言模型多了一个“条件”——多给它一个 $\vect{x}$，它再预测下一个词。Seq2seq 的 decoder
就是这样一个条件语言模型，条件是源序列。
\end{intuitionblock}

\subsubsection{从 LM 到 CLM 的形式化}

普通语言模型把一句话的概率拆成左到右的链式乘积：
\[
    p(\vect{y}) \;=\; \prod_{t=1}^{T} p(y_t \mid y_{<t}),
    \qquad y_{<t} := (y_1, \dots, y_{t-1}).
\]
Seq2seq 模型做的事完全一样，只是\emph{每一步都额外条件}于源序列 $\vect{x}$：
\begin{mathbox}{条件语言模型 (CLM) 的核心公式}
\[
    \boxed{\;
    p(\vect{y} \mid \vect{x}) \;=\; \prod_{t=1}^{T} p(y_t \mid y_{<t},\, \vect{x})
    \;}
\]
\end{mathbox}

也就是说：\emph{seq2seq = 条件语言模型}。它和普通 LM 在结构上几乎一样，差别只在“写每个词的时候，
模型可以额外看一眼 $\vect{x}$”。

\subsubsection{三步流水线}

不论 encoder/decoder 内部用什么神经网络，每一步生成 $y_t$ 都是同样三步：
\begin{enumerate}
    \item \textbf{喂入}：把源序列 $\vect{x}$ 和已生成的前缀 $y_{<t}$ 一起送入网络。
    \item \textbf{得到表示}：网络给出一个 $d$ 维向量 $\vect{h}_t \in \mathbb{R}^d$，
          代表“此刻应该写什么样的词”。
    \item \textbf{打分并归一化}：经过一个线性层 $W \in \mathbb{R}^{|V|\times d}$ 把
          $\vect{h}_t$ 映射到词表 $V$ 上，再用 softmax 得到下一个词的分布：
          \[
              p(y_t \mid y_{<t}, \vect{x})
              \;=\; \softmax(W \vect{h}_t)_{y_t}.
          \]
\end{enumerate}

\begin{remark}
$\vect{x}$ 不必是文本。把图像先用 CNN 编码成向量，再让 decoder 写出一段描述，就是
\emph{image captioning}。所以 CLM 比“文本到文本的 seq2seq”更广泛——它是一种通用模板。
\end{remark}

\begin{engbox}
代码骨架（伪代码）：
\begin{lstlisting}
h_src = encoder(x)                  # 源序列的表示
y_prev = [BOS]
for t in range(T_max):
    h_t   = decoder(y_prev, h_src)  # 第 t 步的隐表示
    logits = W @ h_t                # 线性投影到 |V|
    probs  = softmax(logits)        # 下一个词的分布
    y_t    = sample_or_argmax(probs)
    y_prev.append(y_t)
    if y_t == EOS: break
\end{lstlisting}
\end{engbox}

\subsection{最简单的模型：两个 RNN/LSTM}
\label{sec:two-rnns}

\begin{intuitionblock}
既然要“读完 $\vect{x}$ 再写 $\vect{y}$”，最朴素的实现就是：
用一个 RNN 把源序列从左到右扫一遍，把它的\emph{最后一个隐藏状态}当成“整句话的总结”，
再把这个状态作为另一个 RNN 的初始状态，让它一个词一个词写出目标序列。
\end{intuitionblock}

\subsubsection{形式化}

设 encoder 的隐状态序列为 $\vect{h}_1^{\text{enc}}, \dots, \vect{h}_S^{\text{enc}}$，
decoder 的隐状态序列为 $\vect{h}_1^{\text{dec}}, \dots, \vect{h}_T^{\text{dec}}$。
最朴素的 seq2seq 模型只用一根细线把两者连起来：
\begin{align}
    \vect{h}_s^{\text{enc}} &= \mathrm{RNN}_{\text{enc}}(\vect{h}_{s-1}^{\text{enc}},\, x_s),
        & s &= 1, \dots, S, \\
    \vect{h}_0^{\text{dec}} &= \vect{h}_S^{\text{enc}}, \label{eq:bottleneck}\\
    \vect{h}_t^{\text{dec}} &= \mathrm{RNN}_{\text{dec}}(\vect{h}_{t-1}^{\text{dec}},\, y_{t-1}),
        & t &= 1, \dots, T.
\end{align}
关键就在式 \eqref{eq:bottleneck}：源句的全部信息——长度 $S$ 任意——都被压成
\emph{一个固定维度的向量} $\vect{h}_S^{\text{enc}}$ 传给 decoder。

\begin{keypointblock}
这一根细线就是\emph{固定向量瓶颈} (fixed-size bottleneck)。它在短句上工作良好，
但句子越长，单个向量越塞不下，翻译质量会随句长显著下降。\textbf{这正是第二章 attention 要解决的问题。}
\end{keypointblock}

\subsubsection{两个值得记住的历史细节}

\begin{remark}[Sutskever 等人 2014]
最早把 “两个 LSTM 堆起来做翻译” 这套打通的工作是 Sutskever, Vinyals, Le 在 2014 年的
\emph{Sequence to Sequence Learning with Neural Networks}。它在 WMT 英法翻译上首次让
纯神经方法逼近当时主流的统计机器翻译系统。
\end{remark}

\begin{remark}[反转源句的小技巧]
他们还报告了一个反直觉但很有效的 trick：把源句\emph{倒着}喂给 encoder。这样源句开头的词
（往往对应目标句开头的词）和 decoder 第一步在时间上更接近，梯度路径更短，优化更顺利。
这个技巧本身在 attention 出现后基本被遗忘，但它揭示了瓶颈结构对“距离”非常敏感。
\end{remark}

\subsection{训练：交叉熵损失 (Cross-Entropy Loss)}
\label{sec:ce-loss}

\begin{problemblock}
模型每一步吐出来的是一个长度为 $|V|$ 的概率分布 $p(y_t \mid y_{<t}, \vect{x})$。
而真实的下一个词 $y_t^*$ 只是\emph{词表上的某一个}。怎么把“分布”和“某个词”比出一个损失？
\end{problemblock}

\subsubsection{每一步：把真值看作 one-hot}

把真实的下一个词 $y_t^*$ 写成 one-hot 向量 $\vect{e}_{y_t^*} \in \{0,1\}^{|V|}$：
只有第 $y_t^*$ 个分量是 1，其余为 0。模型的预测分布记为
$\vect{p}_t \in \Delta^{|V|-1}$。两者之间的\emph{交叉熵}定义为
\[
    H(\vect{e}_{y_t^*},\, \vect{p}_t)
    \;=\; -\sum_{w \in V} \vect{e}_{y_t^*}[w] \,\log \vect{p}_t[w]
    \;=\; -\log p(y_t^* \mid y_{<t}, \vect{x}).
\]
因为 one-hot 只有一个位置非零，整个求和\emph{坍缩成单项}：就是模型给真值那个词的负对数概率。

\subsubsection{整个序列：把每步加起来}

\begin{mathbox}{Seq2seq 的训练损失}
\[
    \mathcal{L}(\vect{x}, \vect{y}^*)
    \;=\; -\sum_{t=1}^{T} \log p(y_t^* \mid y_{<t}^*,\, \vect{x}).
\]
在一个 mini-batch 上，对所有句子再取平均（或先除以各自长度后再取平均）。
\end{mathbox}

\subsubsection{为什么是这个损失？}

\begin{theorem}[交叉熵 = 最大似然 = 最小 KL]
最小化上面这个交叉熵损失，等价于在训练数据上做\emph{最大似然估计}：
\[
    \min_{\theta}\; -\sum_{(\vect{x},\vect{y}^*)} \log p_\theta(\vect{y}^* \mid \vect{x})
    \;\Longleftrightarrow\;
    \max_{\theta}\; \prod_{(\vect{x},\vect{y}^*)} p_\theta(\vect{y}^* \mid \vect{x}),
\]
也等价于让模型分布 $p_\theta(\cdot \mid \vect{x})$ 在 KL 意义下尽量\emph{逼近}数据经验分布
$\hat{p}_{\text{data}}(\cdot \mid \vect{x})$。
\end{theorem}

\subsubsection{Teacher Forcing 与暴露偏差}

训练时还有一个\emph{常常被新手忽略}的细节：在算第 $t$ 步的损失时，decoder 的输入“上一个词”
是用真值 $y_{t-1}^*$，\emph{不是}模型自己上一步预测出来的词。这种做法叫做
\textbf{teacher forcing}（教师强迫）。

\begin{itemize}
    \item \textbf{为什么这么做？} 训练初期模型预测几乎全错，如果一直拿自己的错答案当输入，
          后面所有时间步都建立在错误前缀上，损失曲面非常崎岖、根本学不动。用真值前缀可以让
          每一步都在“合理”的上下文里学习，训练快很多。
    \item \textbf{副作用：暴露偏差 (exposure bias)。} 推理时没有真值可用，模型只能看自己之前
          的预测；而训练时它从来没见过自己的错误前缀。这种\emph{训练-推理不一致}会让长序列
          生成时错误累积。后续课程里 scheduled sampling、minimum risk training、RL 微调
          等都是为了缓解这件事，但本章先放一放。
\end{itemize}

\begin{engbox}
训练循环骨架（伪代码）：
\begin{lstlisting}
h_src = encoder(x)
h_dec = h_src                          # 瓶颈传话
loss = 0
for t in range(1, T+1):
    h_dec = decoder_step(h_dec, y_star[t-1])   # teacher forcing
    logits = W @ h_dec
    loss  += cross_entropy(logits, y_star[t])  # = -log p(y*_t | ...)
loss.backward()
\end{lstlisting}
\end{engbox}

\subsection{推理：贪心解码与 Beam Search}
\label{sec:decoding}

\begin{problemblock}
训练好之后，给定一个新的 $\vect{x}$，我们想要的是\emph{最可能}的目标序列：
\[
    \vect{y}^{\star} \;=\; \arg\max_{\vect{y}} \; p(\vect{y} \mid \vect{x}).
\]
但词表大小 $|V|$ 通常上万，目标长度可达几十——可能的 $\vect{y}$ 有 $|V|^T$ 种，
\emph{穷举绝对不可行}。怎么办？
\end{problemblock}

\subsubsection{方案 A：贪心解码 (Greedy Decoding)}

\begin{intuitionblock}
最朴素的近似：每一步都挑当前\emph{条件概率最大}的那个词，写下来，作为下一步的输入，反复直到 EOS。
\end{intuitionblock}
\[
    \hat{y}_t \;=\; \arg\max_{w \in V} \; p(w \mid \hat{y}_{<t},\, \vect{x}).
\]

优点是快、实现极简单；缺点是\emph{局部最优不等于全局最优}。某一步选了概率第一的词，
可能让后续所有步骤都陷入低概率区域；如果当时选了概率第二的词，整句话反而更可能。

\begin{example}
英译中：源句 \emph{``I saw her duck.''}。Greedy 在第三步可能锁定“鸭子”，导致整句翻成
“我看见她的鸭子”；但若回退一步把 \emph{duck} 当动词，整句应是“我看见她躲开”——
后者作为整句的概率更高，却被 greedy 永远错过。
\end{example}

\subsubsection{方案 B：Beam Search}

\begin{intuitionblock}
不要在每一步只留一个候选，而是\emph{始终保留 $N$ 个最有希望的部分序列}（叫 beam）。
每一步把这 $N$ 个候选各扩展 $|V|$ 种延续，得到 $N\cdot|V|$ 个新候选，
\emph{再}按累计 log 概率挑出前 $N$ 个，进入下一步。
\end{intuitionblock}

\begin{mathbox}{Beam Search 的打分}
对一条部分假设 $\vect{y}_{1:t}$，累计得分是 log 概率之和：
\[
    \score(\vect{y}_{1:t} \mid \vect{x})
    \;=\; \sum_{i=1}^{t} \log p(y_i \mid y_{<i},\, \vect{x}).
\]
每步保留 $\score$ 最高的 $N$ 条。$N$ 称为 \emph{beam size}（束宽），机器翻译里常取
$N = 4 \sim 10$。
\end{mathbox}

\begin{itemize}
    \item $N = 1$ 时退化为 greedy。
    \item $N \to \infty$ 时趋近于全局最优搜索（穷举）。
\end{itemize}

\subsubsection{Beam 越大就越好吗？——并不}

经验上，beam size 不是越大越好：太大反而会让翻译质量\emph{下降}。常见的解释有两类：

\begin{itemize}
    \item \textbf{长度偏置 (length bias)。} log 概率是负数，每加一个词都让总分变小，
          所以 beam 越大、越倾向于挑\emph{更短}的句子，可能会丢掉应有的内容。
          常用的修补是\emph{长度归一化} (length normalization)：用
          $\score(\vect{y}_{1:t}) / t^{\alpha}$（$\alpha \in [0.6, 1]$）作为打分，
          抵消长度对累计 log 概率的拉低。
    \item \textbf{模型本身的过自信。} 训练时只见过真实序列，模型把概率质量过度集中在某些
          “典型”句子上；当 beam 足够大时，搜索算法会“恰好”找到一些极短或退化的高分
          序列（例如直接输出 EOS）。这是模型问题，不是搜索问题。
\end{itemize}

\begin{engbox}
Beam search 的最小骨架（伪代码）：
\begin{lstlisting}
beams = [(["BOS"], 0.0)]               # (前缀, 累计 log 概率)
for t in range(T_max):
    candidates = []
    for prefix, score in beams:
        if prefix[-1] == "EOS":
            candidates.append((prefix, score))
            continue
        logp = log_softmax(decoder_step(prefix, x))   # 形状 (|V|,)
        for w, lp in topk(logp, k=N):                 # 只看每条的 top-N 延续
            candidates.append((prefix + [w], score + lp))
    beams = topk_by_score(candidates, k=N)            # 全局再挑前 N
return best_by_length_normalized_score(beams)
\end{lstlisting}
\end{engbox}

\subsection{这一章真正该记住的}

\begin{keypointblock}
\begin{enumerate}
    \item \textbf{Seq2seq = encoder + decoder。} 一边读源序列、一边写目标序列；后续所有更复杂的
          模型都没有抛弃这个骨架，只是在改两边怎么实现、怎么连接。
    \item \textbf{Decoder 是一个条件语言模型}：
          $p(\vect{y} \mid \vect{x}) = \prod_t p(y_t \mid y_{<t}, \vect{x})$。
          它和普通 LM 唯一的区别就是\emph{额外条件于 $\vect{x}$}。
    \item \textbf{最朴素实现 = 两个 RNN + 一个固定向量瓶颈。} 这个瓶颈对长句不友好，
          正是第二章 attention 要解决的痛点。
    \item \textbf{训练用 token 级交叉熵}，等价于对训练序列做最大似然。Decoder 输入用真值前缀
          (\emph{teacher forcing})，会带来训练-推理不一致 (\emph{exposure bias})。
    \item \textbf{推理是搜索 $\arg\max_{\vect{y}} p(\vect{y} \mid \vect{x})$ 的近似。}
          Greedy 局部最优、Beam search 用宽度为 $N$ 的近似全局搜索，常用 $N \approx 4\text{--}10$。
    \item \textbf{Beam 不是越大越好}，原因是长度偏置和模型过自信；
          长度归一化是最常见的实用修补。
\end{enumerate}
\end{keypointblock}

\section{Attention：让解码器自己去原文里取信息}
\label{ch:attention}

\subsection{章节定位 + 零基础该问的问题}

上一章我们看到，经典的 seq2seq 把整句源语言压成一个向量 $\vect{h}^{\mathrm{enc}}_m$，再交给解码器。这一章要回答的问题是：当这个“压缩到一个向量”的方案吃力的时候，模型应该怎么改？答案就是 \emph{attention}。它是 Transformer 之前 NMT 时代最重要的进展，也是 self-attention 的直接祖先。

\begin{problemblock}[一上来该问的几个问题]
带着这些问题往下读，比一上来就背公式更划算：
\begin{enumerate}
  \item RNN 编码器到底\emph{卡在哪}？为什么句子一长就翻不动？
  \item Attention 的输入是什么、输出是什么？它放在 seq2seq 的哪个位置？
  \item Attention score 函数有几种，互相差在哪？
  \item Bahdanau attention 和 Luong attention 是同一个东西吗？哪里不一样？
  \item 学出来的 attention 权重，为什么大家都说它“像对齐”？这有什么用？
\end{enumerate}
\end{problemblock}

\subsection{固定编码器表示的瓶颈}
\label{sec:bottleneck}

\subsubsection{背景问题：一个向量装不下整句话}

\begin{problemblock}
经典 seq2seq 的编码器读完全句之后，吐出一个固定维度的向量 $\vect{h}^{\mathrm{enc}}_m \in \R^{d}$，把它当作整句的“意思”交给解码器。但是：
\begin{itemize}
  \item 句子是 5 个词还是 50 个词，最后给解码器的都是同一个 $d$ 维向量。
  \item 解码器在生成第 1 个目标词时关心的源信息，和生成第 10 个目标词时关心的源信息，\emph{显然不一样}，但它每一步看到的“源信息”都是同一个向量。
\end{itemize}
\end{problemblock}

\subsubsection{朴素直觉：把宇宙压进 512 个数}

\begin{intuitionblock}
Lena 在博客里有一个非常形象的比喻：让编码器把整句话压成一个固定向量，就像让你\emph{把整个宇宙压进 512 个数}里再交出去。短句还行，长句就一定会丢信息。

更糟的是：解码器每一步都要从同一个“宇宙摘要”里翻出当前需要的那一小块信息——它根本没机会“回头再看一眼原文”。
\end{intuitionblock}

\subsubsection{数学表达：瓶颈到底在哪一项}

\begin{mathbox}{固定编码器表示的损失}
经典 seq2seq 的解码概率分解为
\begin{equation}
p(\vect{y}\mid \vect{x}) \;=\; \prod_{t=1}^{n} p\!\left(y_t \mid y_{<t},\,\vect{h}^{\mathrm{enc}}_m\right).
\end{equation}
注意右边条件里的 $\vect{h}^{\mathrm{enc}}_m$ 与 $t$ 无关。也就是说，模型被强制要求：
\begin{equation}
\text{对所有解码步 } t,\ \text{源端信息都用\emph{同一个}向量 }\vect{h}^{\mathrm{enc}}_m\text{ 表达。}
\end{equation}
当 $m$ 增大时，这个向量必须承担越来越多内容，信息被覆盖、被冲掉。
\end{mathbox}

\begin{remark}
经验上，普通 seq2seq 在源句长度超过 20--30 之后 BLEU 会明显下降；引入 attention 之后这一掉点会被显著拉平。这就是 Bahdanau 等人 2014 年那篇论文最早展示的结果。
\end{remark}

\subsection{Attention 的高层视角}
\label{sec:attn-overview}

\subsubsection{背景问题：能不能让解码器“回头看”？}

\begin{problemblock}
既然“压成一个向量”是瓶颈，那很自然的想法是：
\begin{itemize}
  \item 让编码器把每个时刻的隐状态 $\vect{s}_1,\ldots,\vect{s}_m$ 都\emph{留下来}，不要扔。
  \item 解码器在第 $t$ 步，根据当前自己的状态 $\vect{h}_t$，自己决定该\emph{重点看}源端的哪些位置。
\end{itemize}
关键约束是：这个“决定看哪里”的过程必须\emph{可微}，否则没法用反向传播一起训练。
\end{problemblock}

\subsubsection{朴素直觉：加权查表}

\begin{intuitionblock}
把编码器输出 $\vect{s}_1,\ldots,\vect{s}_m$ 想成一张“源端记忆表”。解码器每一步做的事就是一次\emph{软查表}：
\begin{enumerate}
  \item 拿出一个“查询” $\vect{h}_t$（解码器当前状态）。
  \item 对每条记忆 $\vect{s}_k$ 打一个相关性分数 $e_{tk}$。
  \item 把分数过一个 softmax 得到“看每条记忆的比例” $\alpha_{tk}$。
  \item 按比例把所有记忆加权求和，得到这一步的“context” $\vect{c}_t$。
\end{enumerate}
注意整个过程没有 argmax、没有硬选择，所以\emph{处处可导}，可以端到端训练。模型自己会学出“在第 $t$ 步该看哪里”，不需要我们提供人工对齐。
\end{intuitionblock}

\subsubsection{数学表达：attention 的三步公式}

\begin{mathbox}{解码第 $t$ 步的 attention}
设编码器输出为 $\vect{s}_1,\ldots,\vect{s}_m \in \R^{d_s}$，解码器当前状态为 $\vect{h}_t \in \R^{d_h}$。
\begin{align}
e_{tk} &\;=\; \score(\vect{h}_t, \vect{s}_k), \qquad k=1,\ldots,m, \label{eq:attn-score}\\
\alpha_{t\cdot} &\;=\; \softmax(e_{t\cdot}) \;\Longleftrightarrow\; \alpha_{tk} \;=\; \frac{\exp(e_{tk})}{\sum_{j=1}^{m}\exp(e_{tj})}, \label{eq:attn-softmax}\\
\vect{c}_t &\;=\; \sum_{k=1}^{m} \alpha_{tk}\,\vect{s}_k. \label{eq:attn-context}
\end{align}
得到的 $\vect{c}_t$ 就是“给第 $t$ 步用的源端摘要”，它会被喂回解码器（拼接到输入、或拼接到隐状态、或参与最终输出层）。
\end{mathbox}

\begin{definition}[Attention 权重]
向量 $\boldsymbol{\alpha}_{t\cdot}=(\alpha_{t1},\ldots,\alpha_{tm})$ 满足 $\alpha_{tk}\ge 0$ 且 $\sum_k \alpha_{tk}=1$，称为解码第 $t$ 步对源端的 \emph{attention 权重}。它描述了“第 $t$ 步用了源端哪些位置、各占多少”。
\end{definition}

\begin{engbox}
工程上要记住三件事：
\begin{itemize}
  \item 编码器需要把\emph{所有时刻}的 $\vect{s}_k$ 都缓存下来（一个 $m\times d_s$ 矩阵），而不是只留最后一个。
  \item 每个解码步要做一次 $O(m)$ 的打分 + softmax + 加权和；总成本相对原 seq2seq 增加了一项 $O(n m)$。
  \item $\vect{c}_t$ 怎么“喂回”解码器，是 Bahdanau 与 Luong 两派的主要差异（见 \S\ref{sec:bahdanau-luong}）。
\end{itemize}
\end{engbox}

\subsection{怎么算 attention score？}
\label{sec:score}

\subsubsection{背景问题：一个 score 函数能管所有情况吗}

\begin{problemblock}
公式 \eqref{eq:attn-score} 里的 $\score(\vect{h}_t,\vect{s}_k)$ 只是一个抽象记号。它要满足两点：
\begin{itemize}
  \item 输入两个向量，输出一个标量分数；
  \item 是 $\vect{h}_t$ 和 $\vect{s}_k$ 的可微函数。
\end{itemize}
满足这两点的函数有一堆，历史上常用的就是下面三种。
\end{problemblock}

\subsubsection{朴素直觉：从“纯几何”到“多塞一层 MLP”}

\begin{intuitionblock}
\begin{itemize}
  \item \textbf{点积}：直接看两个向量的余弦相似度（按模长缩放就成了后来的 scaled dot-product）。最便宜，但要求 $\vect{h}_t$ 和 $\vect{s}_k$ 在\emph{同一空间}里有意义。
  \item \textbf{双线性 (bilinear)}：在中间塞一个矩阵 $\vect{W}$，相当于先把其中一个向量投影到对方的空间。$\vect{h}_t$ 和 $\vect{s}_k$ 维度不同也能用。
  \item \textbf{加性 (additive / MLP)}：把两个向量先各自线性变换、加起来，再过 $\tanh$，再投到一个标量。最“能学”，但参数和算力都最贵。
\end{itemize}
\end{intuitionblock}

\subsubsection{数学表达：三种常见 score 函数}

\begin{mathbox}{三种 attention score}
\begin{align}
\textbf{点积 (dot)}: \quad &\score(\vect{h},\vect{s}) \;=\; \vect{h}^{\top}\vect{s}, \label{eq:score-dot}\\
\textbf{双线性 (Luong)}: \quad &\score(\vect{h},\vect{s}) \;=\; \vect{h}^{\top}\vect{W}\vect{s}, \label{eq:score-bilinear}\\
\textbf{加性 (Bahdanau)}: \quad &\score(\vect{h},\vect{s}) \;=\; \vect{v}^{\top}\tanh\!\left(\vect{W}_1 \vect{h} + \vect{W}_2 \vect{s}\right). \label{eq:score-additive}
\end{align}
其中 $\vect{W}\in\R^{d_h\times d_s}$，$\vect{W}_1\in\R^{d\times d_h}$，$\vect{W}_2\in\R^{d\times d_s}$，$\vect{v}\in\R^{d}$ 都是可学习参数。
\end{mathbox}

\begin{remark}[谁更好用]
\begin{itemize}
  \item \textbf{速度}：点积 $\gg$ 双线性 $>$ 加性。点积可以一次矩阵乘法批量算完，所以后来 Transformer 选了它。
  \item \textbf{表达力}：加性最强（多了一层非线性），点积最弱。
  \item \textbf{兼容性}：维度不一致时只能用双线性或加性。
\end{itemize}
\end{remark}

\subsection{模型变体：Bahdanau vs Luong}
\label{sec:bahdanau-luong}

\subsubsection{背景问题：attention 应该插在解码器的哪一步}

\begin{problemblock}
有了 attention 模块以后，下一个问题是：
\begin{itemize}
  \item 解码器更新隐状态的\emph{之前}就先算 $\vect{c}_t$，把它和输入一起喂给 RNN？还是
  \item 让 RNN 先正常更新出 $\vect{h}_t$，再用 $\vect{h}_t$ 去算 $\vect{c}_t$，最后把两者合起来出预测？
\end{itemize}
这就是 Bahdanau (2014) 和 Luong (2015) 两派的核心区别。
\end{problemblock}

\subsubsection{朴素直觉：时序对比}

\begin{intuitionblock}
两条时序线，一眼对比：
\begin{itemize}
  \item \textbf{Bahdanau}：$\vect{h}_{t-1}\ \longrightarrow\ \text{attention}\ \longrightarrow\ \vect{c}_t\ \longrightarrow\ \text{RNN 用 }(\vect{c}_t, y_{t-1})\text{ 算 }\vect{h}_t\ \longrightarrow\ \text{预测 }y_t$。\\ Attention 在“解码 RNN 一步”\emph{之前}发生。
  \item \textbf{Luong}：$\vect{h}_{t-1}\ \longrightarrow\ \text{RNN 用 }y_{t-1}\text{ 算 }\vect{h}_t\ \longrightarrow\ \text{attention(}\vect{h}_t\text{)}\ \longrightarrow\ \vect{c}_t\ \longrightarrow\ \tilde{\vect{h}}_t = \tanh(\vect{W}[\vect{h}_t;\vect{c}_t])\ \longrightarrow\ \text{预测 }y_t$。\\ Attention 在“解码 RNN 一步”\emph{之后}发生。
\end{itemize}
\end{intuitionblock}

\subsubsection{数学表达：两个版本的关键公式}

\begin{mathbox}{Bahdanau attention（2014）}
\begin{itemize}
  \item 编码器是\emph{双向} RNN：$\vect{s}_k = [\overrightarrow{\vect{s}}_k;\overleftarrow{\vect{s}}_k]$。
  \item Score 用加性：$e_{tk} = \vect{v}^{\top}\tanh(\vect{W}_1 \vect{h}_{t-1} + \vect{W}_2 \vect{s}_k)$。
  \item 注意 query 是\emph{上一步}的 $\vect{h}_{t-1}$，先算 $\vect{c}_t$，再更新：
\end{itemize}
\begin{align}
e_{tk} &= \vect{v}^{\top}\tanh(\vect{W}_1 \vect{h}_{t-1} + \vect{W}_2 \vect{s}_k),\\
\alpha_{t\cdot} &= \softmax(e_{t\cdot}),\quad \vect{c}_t = \sum_k \alpha_{tk}\vect{s}_k,\\
\vect{h}_t &= \mathrm{RNN}\!\left(\vect{h}_{t-1},\ [\vect{E}\, y_{t-1};\,\vect{c}_t]\right),\\
p(y_t\mid y_{<t},\vect{x}) &= \softmax\!\left(\vect{W}_o[\vect{h}_t;\vect{c}_t;\vect{E}\,y_{t-1}]\right).
\end{align}
\end{mathbox}

\begin{mathbox}{Luong attention（2015，global 版本）}
\begin{itemize}
  \item 编码器是\emph{单向} RNN（更简单）。
  \item Score 用双线性：$e_{tk} = \vect{h}_t^{\top}\vect{W}\vect{s}_k$。
  \item Query 是\emph{当前步}的 $\vect{h}_t$，先更新 RNN，再算 attention：
\end{itemize}
\begin{align}
\vect{h}_t &= \mathrm{RNN}\!\left(\vect{h}_{t-1},\ \vect{E}\, y_{t-1}\right),\\
e_{tk} &= \vect{h}_t^{\top}\vect{W}\vect{s}_k,\quad \alpha_{t\cdot}=\softmax(e_{t\cdot}),\quad \vect{c}_t=\sum_k \alpha_{tk}\vect{s}_k,\\
\tilde{\vect{h}}_t &= \tanh\!\left(\vect{W}_c[\vect{h}_t;\vect{c}_t]\right),\\
p(y_t\mid y_{<t},\vect{x}) &= \softmax\!\left(\vect{W}_o\,\tilde{\vect{h}}_t\right).
\end{align}
\end{mathbox}

\begin{remark}
Luong 还提出过一个 \emph{local attention} 版本，只在源端附近一个窗口里算 softmax。但工业界后来还是 global 版本+各种改进用得多。
\end{remark}

\begin{engbox}
理解这两派的差异比记公式重要：
\begin{itemize}
  \item \textbf{时序}：Bahdanau 是“先看再走”，Luong 是“先走再看”。
  \item \textbf{Score}：Bahdanau 用 MLP（参数最多、最能学）；Luong 默认用双线性（更轻）。
  \item \textbf{编码器}：Bahdanau 双向（信息更全），Luong 单向（实现更简单）。
  \item 这两个模型在 Transformer 出现前\emph{统治了}神经机器翻译，几乎所有后续 NMT 工程都是它们的变体。
\end{itemize}
\end{engbox}

\subsection{Attention 学到了（近似的）对齐}
\label{sec:alignment}

\subsubsection{背景问题：这些 $\alpha_{tk}$ 真的有意义吗}

\begin{problemblock}
我们没有给模型任何“哪个目标词对应哪个源词”的标注，只是用最大似然让它去预测下一个目标词。可学完之后，把 $\alpha_{tk}$ 画成热力图，往往会看到：解码到目标位置 $t$ 时，权重恰好集中在源端那几个“正在被翻译”的词上。
\end{problemblock}

\subsubsection{朴素直觉：软对齐}

\begin{intuitionblock}
统计机器翻译时代，“对齐 (alignment)”是一个核心概念，要靠 IBM Model 1/2 等专门的算法去估。Attention 把这件事顺手做了：
\begin{itemize}
  \item 每个目标词 $y_t$ 对应一个源端的概率分布 $\boldsymbol{\alpha}_{t\cdot}$。
  \item 这就是“软对齐”——不再是 hard 的“词 $i$ 对应词 $j$”，而是“词 $t$ 用了源端 30\% 的位置 $j_1$，60\% 的位置 $j_2$，\ldots”。
\end{itemize}
\end{intuitionblock}

\subsubsection{从可解释性的角度看}

\begin{remark}[Attention 作为归纳偏置]
Bahdanau 那篇论文里展示的 attention 热力图，是深度学习历史上\emph{第一批}看起来很“可解释”的内部状态：
\begin{itemize}
  \item 它没有被强行训练成对齐，但自发地呈现出对齐的样子；
  \item 这说明“对每一步只看相关的几条记忆”是一种\emph{有用}且\emph{自然}的归纳偏置；
  \item 这一思想被 Transformer 推到极致：\emph{所有}信息流都通过 attention 完成。
\end{itemize}
不过要小心：后续研究（如 Jain \& Wallace 2019）也指出，attention 权重并不一定等于“真正的解释”。它\emph{看起来}像对齐，并不代表模型一定\emph{用}它做决策。
\end{remark}

\subsection{这一章真正该记住的}

\begin{keypointblock}
\begin{enumerate}
  \item \textbf{病根}：经典 seq2seq 把整句源语言压成\emph{一个固定向量}；句子一长，就装不下。
  \item \textbf{治法}：编码器保留所有时刻的状态 $\vect{s}_1,\ldots,\vect{s}_m$；解码器每一步对它们做\emph{加权查表}，权重由 softmax 给出。
  \item \textbf{三步范式}：score $\to$ softmax $\to$ 加权和；这套流程全程可导，端到端训练，不需要人工对齐。
  \item \textbf{Score 三选一}：点积（最快）、双线性（Luong）、加性 MLP（Bahdanau）；速度 vs 表达力的取舍。
  \item \textbf{Bahdanau vs Luong}：核心差异是“attention 在 RNN 更新之前还是之后”，以及 score 函数和编码器方向；它们是 Transformer 之前 NMT 的两大支柱。
  \item \textbf{学出来像对齐}：$\alpha_{tk}$ 自发呈现软对齐，是“attention 作为归纳偏置”的第一个直观证据；它直接为下一章的 self-attention 与 Transformer 铺好了路。
\end{enumerate}
\end{keypointblock}

\section{Transformer：Attention is All You Need}
\label{sec:ch03-transformer}

\subsection{章节定位}
\label{subsec:ch03-positioning}

\begin{problemblock}
读完第二章，我们已经知道 attention 可以让 decoder 在每一步“挑”encoder 的不同位置看。
但 RNN 仍然是骨架：编码端要顺序读完整个源句，解码端要顺序生成。这意味着\textbf{长序列训练慢、难以并行}，并且远距离依赖只能通过隐状态间接传递。

那么自然要问：能不能完全不用 RNN，只用 attention 就把序列建模做完？
\end{problemblock}

\begin{intuitionblock}
2017 年 Vaswani 等人的论文《Attention Is All You Need》给出的回答是：可以。
他们提出的 \textbf{Transformer} 架构\textbf{不含任何 recurrence、也不含卷积}，整套模型只由 attention、线性层、归一化和位置编码拼成。

它带来的直接好处：
\begin{itemize}
  \item \textbf{训练阶段并行}：一句话里所有位置可以同时算，而不是按时间步递推。
  \item \textbf{长依赖一步到位}：任意两个位置之间的“距离”都是 1 次 attention，而不是 $O(n)$ 步隐状态传递。
  \item \textbf{已成为事实标准}：今天几乎所有大语言模型 (GPT/Llama/Qwen/\ldots) 都是 Transformer 的变种。
\end{itemize}
\end{intuitionblock}

为了把这一章读“懂”而不是只“看过”，我们围绕 5 个新手最容易卡住的问题来展开：
\begin{enumerate}
  \item 什么叫\textbf{自注意力 (self-attention)}？跟第二章那个 encoder--decoder 之间的 attention 有什么区别？
  \item 为什么一个 token 要同时扮演 \textbf{Query / Key / Value} 三个角色，不能只用一个向量？
  \item 公式里那个 $\sqrt{d_k}$ 是哪来的，能不能不除？
  \item \textbf{Masked} self-attention 是在 mask 谁、为什么要 mask？
  \item \textbf{Multi-head} 到底比单头多做了什么？
\end{enumerate}
后面每一节会回答其中一两个。

\subsection{自注意力 Self-Attention：让 token 互相看}
\label{subsec:ch03-selfattn}

\begin{problemblock}
第二章的 attention 是\textbf{跨序列}的：decoder 的当前位置当 query，去看 encoder 的所有位置。
但编码端自己内部呢？“The animal didn't cross the street because \emph{it} was too tired.” 里面 ``it'' 到底指 ``animal'' 还是 ``street''？
我们希望 encoder 的同一层里，每个词都能看一眼同一句话里的其他词，再决定自己的表示。
\end{problemblock}

\begin{intuitionblock}
\textbf{Self-attention} 就是把 query、key、value \emph{都来自同一侧} 的 attention。
对一句话里的每个 token，做这样三件事：
\begin{itemize}
  \item 我把自己变成一个 \textbf{Query} 向量 $\vect{q}$：表示\emph{“我想问什么”}。
  \item 我把自己变成一个 \textbf{Key} 向量 $\vect{k}$：表示\emph{“别人来找我时，我对外的标签是什么”}。
  \item 我把自己变成一个 \textbf{Value} 向量 $\vect{v}$：表示\emph{“如果别人决定关注我，我会贡献什么内容”}。
\end{itemize}
然后让每个 query 跟所有 key 比相似度，得到权重，再用这些权重对所有 value 加权求和，
就得到这个 token 的新表示。整句话里所有 token \emph{并行}做这件事。
\end{intuitionblock}

数学上，把整句话堆成矩阵：$Q \in \R^{n \times d_k}$，$K \in \R^{n \times d_k}$，$V \in \R^{n \times d_v}$，其中 $n$ 是序列长度。

\begin{mathbox}{Scaled Dot-Product Attention}
\[
  \mathrm{Attention}(Q, K, V)
  \;=\; \softmax\!\left(\frac{Q K^\top}{\sqrt{d_k}}\right) V .
\]
其中：
\begin{itemize}
  \item $Q K^\top \in \R^{n \times n}$：第 $(i,j)$ 项就是位置 $i$ 的 query 和位置 $j$ 的 key 的内积，
        相当于 $\score(\vect{q}_i, \vect{k}_j)$。
  \item 除以 $\sqrt{d_k}$：缩放因子，下面解释。
  \item $\softmax$ 沿\textbf{每一行}做：每个 query 对所有 key 的权重和为 1。
  \item 最后乘 $V$：用这些权重把 value 加权求和。
\end{itemize}
\end{mathbox}

\begin{remark}[为什么要除 $\sqrt{d_k}$？]
两个独立、零均值、单位方差的 $d_k$ 维向量做点积，结果的方差大约是 $d_k$。
$d_k$ 一大（例如 64、128），点积的绝对值就会很大；进入 softmax 后，
最大那一项会被放得极大，其它几乎归零，softmax 退化成 \emph{近似 hard argmax}。
这时\textbf{梯度几乎为 0}，训练走不动。除以 $\sqrt{d_k}$ 把方差压回 $\approx 1$，softmax 才有梯度。
\end{remark}

\begin{engbox}
\textbf{Q/K/V 是怎么来的？} 不是凭空写的，而是用三个可学习的投影矩阵从同一份输入 $X \in \R^{n \times d_{\text{model}}}$ 算出来：
\[
  Q = X W_Q,\quad K = X W_K,\quad V = X W_V .
\]
也就是说\textbf{“同一个 token 的三种角色”都是它自己的线性变换}，模型自己学三个角度的“画像”。
\end{engbox}

\begin{example}[“it” 该指谁]
继续上面的例子：当 ``it'' 这个位置的 query 跟 ``animal'' 的 key 内积大、跟 ``street'' 的 key 内积小，
softmax 之后 ``animal'' 的权重接近 1，``it'' 的新表示就主要由 ``animal'' 的 value 组成。
模型不需要任何递归，一步就把指代关系拉到了眼前。
\end{example}

\subsection{Masked Self-Attention：decoder 不许偷看未来}
\label{subsec:ch03-masked}

\begin{problemblock}
Decoder 在\textbf{推理时}是逐词生成的：先吐 $y_1$，再根据 $y_1$ 吐 $y_2$\ldots\ 它\emph{物理上}看不到未来。

但\textbf{训练时}不一样：为了并行，我们会把整句目标序列 $y_1, y_2, \dots, y_T$ 一次性喂进 decoder。
如果直接做 self-attention，位置 $t$ 的 query 会看到 $y_{t+1}, y_{t+2}, \dots$ 的 key/value，
等于直接把答案抄过来——训练 loss 会瞬间归零，但模型啥也没学会。
\end{problemblock}

\begin{intuitionblock}
解决办法：在 softmax 之前\textbf{把“未来位置”的分数手动设成 $-\infty$}。
$\mathrm{exp}(-\infty) = 0$，softmax 之后这些位置的权重自然为 0，
就像它们根本不存在一样。这样\textbf{并行训练}和\textbf{推理时只能看过去}就保持一致。
\end{intuitionblock}

\begin{mathbox}{Causal Mask}
设 $S = QK^\top / \sqrt{d_k} \in \R^{n \times n}$ 是未 mask 的分数矩阵，定义因果 mask
\[
  M_{ij} =
  \begin{cases}
    0,        & j \le i \\
    -\infty,  & j > i
  \end{cases}
\]
则 masked self-attention 为
\[
  \mathrm{Attention}_{\text{masked}}(Q,K,V)
  \;=\; \softmax(S + M)\, V .
\]
矩阵 $S + M$ 是一个\textbf{下三角分数矩阵}，softmax 沿每行做之后上三角全部为 0。
\end{mathbox}

\begin{remark}[和 RNN decoder 的对比]
RNN decoder 训练一句长度 $T$ 的目标，需要 \textbf{$T$ 步顺序}前向传播；
masked self-attention 把这 $T$ 步压成\textbf{一次}矩阵乘法，
每个位置同时算自己的输出，但通过 mask 保证它\emph{只用得到自己以前的信息}。
这就是 Transformer 训练比 RNN 快很多的关键之一。
\end{remark}

\subsection{多头注意力 Multi-Head Attention}
\label{subsec:ch03-mha}

\begin{problemblock}
一个词通常跟句子里的别的词存在\textbf{多种不同的关系}：
主谓一致、修饰、指代、跨从句的并列\ldots\
单头 attention 只能学一种 ``关注模式''——softmax 那一行决定大头给谁，剩下的就只能分到很小的权重。
如果我希望同一句话里既能抓 ``it $\to$ animal'' 的指代，又能抓 ``didn't $\to$ cross'' 的否定关系，单头不够用。
\end{problemblock}

\begin{intuitionblock}
开\textbf{多个并行的注意力头}：每个头有自己独立的 $W_Q^i, W_K^i, W_V^i$，
独立地学一种 ``关注什么''。最后把所有头的输出\textbf{拼接}起来，再过一个统一的输出投影 $W_O$ 融合一下。
\end{intuitionblock}

\begin{mathbox}{Multi-Head Attention}
设有 $h$ 个头，每个头的维度为 $d_k = d_v = d_{\text{model}} / h$。第 $i$ 个头：
\[
  \mathrm{head}_i \;=\; \mathrm{Attention}\!\bigl(Q W_Q^i,\; K W_K^i,\; V W_V^i\bigr).
\]
多头注意力的输出为
\[
  \mathrm{MultiHead}(Q, K, V)
  \;=\; \mathrm{Concat}(\mathrm{head}_1, \dots, \mathrm{head}_h)\, W_O .
\]
其中 $W_O \in \R^{d_{\text{model}} \times d_{\text{model}}}$ 是输出投影。
\end{mathbox}

\begin{engbox}
\textbf{重要的实现细节：多头\emph{不会}让模型变大。}
工程上不会真的开 $h$ 套独立的小矩阵，而是：
\begin{enumerate}
  \item 用\textbf{一个}大的线性层一次性算出 $Q, K, V$，每个的形状是 $(B, n, d_{\text{model}})$；
  \item 把最后一维 \textbf{reshape} 成 $(h, d_{\text{model}}/h)$，再 \textbf{transpose} 让 $h$ 变成 batch 轴；
  \item 形状变成 $(B, h, n, d_k)$，然后照常做 scaled dot-product，等价于 $h$ 个头\textbf{并行}计算；
  \item 算完再 transpose 回去 reshape 拼接，过 $W_O$。
\end{enumerate}
所以同样的 $d_{\text{model}}$ 下，``单头'' 和 ``多头'' 的参数量是\textbf{一样的}，只是把同一份参数\emph{切}成几份去看不同关系。
\end{engbox}

下面是一段 PyTorch 风格的伪代码，演示 reshape + transpose 这个技巧：

\begin{lstlisting}[language=Python]
# x: (B, n, d_model)
qkv = self.qkv_proj(x)             # (B, n, 3 * d_model)
q, k, v = qkv.chunk(3, dim=-1)     # each: (B, n, d_model)

# split heads: d_model -> (h, d_k)
def split(t):
    return t.view(B, n, h, d_k).transpose(1, 2)   # (B, h, n, d_k)

q, k, v = split(q), split(k), split(v)

scores = (q @ k.transpose(-2, -1)) / math.sqrt(d_k)   # (B, h, n, n)
if causal_mask is not None:
    scores = scores + causal_mask                      # -inf on j > i
attn = scores.softmax(dim=-1) @ v                      # (B, h, n, d_k)

# merge heads back: (B, h, n, d_k) -> (B, n, d_model)
out = attn.transpose(1, 2).contiguous().view(B, n, d_model)
out = self.o_proj(out)                                 # (B, n, d_model)
\end{lstlisting}

\subsection{Transformer 模型架构}
\label{subsec:ch03-arch}

到这里 attention 这块的零件都齐了。剩下的事是把它们和几个``脚手架''组件拼成一个完整的 encoder/decoder 模型。

\subsubsection{Feed-Forward Block (FFN)}

\begin{problemblock}
Self-attention 的工作是``\textbf{看别的位置、把信息搬过来}''。但搬完之后，这个位置自己还需要一步\textbf{非线性的处理}，
把刚刚收集到的混合信息消化成下一层能用的特征。如果只堆 attention，没有非线性变换，整个模型在线性能力上会很受限。
\end{problemblock}

\begin{intuitionblock}
在每一层 attention 之后，对\textbf{每个位置独立}地过一个 2 层 MLP：
先升维、做 ReLU、再降回原维度。位置之间不交互——因为 ``跨位置交互'' 已经由 attention 负责了，FFN 负责``\textbf{各自消化}''。
\end{intuitionblock}

\begin{mathbox}{Position-wise Feed-Forward}
\[
  \mathrm{FFN}(x) \;=\; \max(0,\; x W_1 + b_1)\, W_2 + b_2 .
\]
通常 $W_1 \in \R^{d_{\text{model}} \times d_{\text{ff}}}$、$W_2 \in \R^{d_{\text{ff}} \times d_{\text{model}}}$，
其中 $d_{\text{ff}}$ 一般取 $4 d_{\text{model}}$。``per-position'' 的意思是：同一个 $W_1, W_2$ 应用到序列里每个 token，token 之间不串。
\end{mathbox}

\subsubsection{残差连接 Residual Connections}

\begin{intuitionblock}
Transformer 把 attention 和 FFN 都包进同一个模板里：
\[
  \text{output} \;=\; \mathrm{LayerNorm}\bigl(x + \mathrm{Sublayer}(x)\bigr) .
\]
这个\textbf{加法}就是``Add \& Norm''里的``Add'',
也就是 ResNet 那种残差连接：让信号有一条``\textbf{直通车}''可以绕过子层。

好处：
\begin{itemize}
  \item 梯度更容易传到底层，可以放心堆几十层。
  \item 子层只需要学``\emph{在原表示上加什么修正}''，比从零学整张表示容易得多。
\end{itemize}
\end{intuitionblock}

每一个 attention 子块和每一个 FFN 子块后面都加这种残差连接。

\subsubsection{层归一化 Layer Normalization}

\begin{problemblock}
深层网络里激活值的尺度容易飘——某一层突然很大或很小，会让训练不稳定。
但 \textbf{BatchNorm} 在变长序列上不太合适：不同样本长度不同，按 batch 统计意义不明确，而且\textbf{推理时一句一句来}，根本没有 batch 统计。
\end{problemblock}

\begin{intuitionblock}
换成 \textbf{LayerNorm}：对\textbf{每一个 token 自己的特征向量}做归一化，
也就是沿着 $d_{\text{model}}$ 这一维算均值和方差，跟 batch 大小、序列长度都无关。
\end{intuitionblock}

\begin{mathbox}{Layer Normalization}
对一个 token 向量 $x \in \R^{d_{\text{model}}}$，
\[
  \mu = \tfrac{1}{d}\sum_{j} x_j,
  \qquad
  \sigma = \sqrt{\tfrac{1}{d}\sum_{j}(x_j - \mu)^2 + \varepsilon},
\]
\[
  \mathrm{LN}(x) \;=\; \frac{x - \mu}{\sigma}\odot \mathit{scale} \;+\; \mathit{bias},
\]
其中 $\mathit{scale}, \mathit{bias} \in \R^{d_{\text{model}}}$ 是可学习参数。
\end{mathbox}

\begin{remark}
LayerNorm 不依赖 batch 维 $\Rightarrow$ 对变长序列友好；训练和推理行为完全一致，不需要像 BN 那样维护 running mean/var。
这也是 NLP 模型几乎清一色用 LayerNorm 而不是 BatchNorm 的原因。
\end{remark}

\subsubsection{位置编码 Positional Encoding}

\begin{problemblock}
Self-attention 的公式 $\softmax(QK^\top/\sqrt{d_k})V$ 里\textbf{没有任何关于位置的信息}：
如果你把输入 token 顺序打乱，输出也只是相应被打乱——模型本身是\textbf{置换等变}的。
但语言显然不是词袋，``狗咬人''和``人咬狗''意思不同，模型必须知道每个 token 在序列中的位置。
\end{problemblock}

\begin{intuitionblock}
最简单的做法：给每个位置 $\mathit{pos}$ 配一个``位置向量'' $\mathrm{PE}_{\mathit{pos}} \in \R^{d_{\text{model}}}$，
直接\textbf{加到 token embedding 上}：
\[
  \tilde{x}_{\mathit{pos}} = \mathrm{Embed}(\text{token}_{\mathit{pos}}) + \mathrm{PE}_{\mathit{pos}} .
\]
然后再送入 self-attention。这样同一个词在不同位置就有不同的输入向量，模型自然能区分顺序。
\end{intuitionblock}

\begin{mathbox}{Sinusoidal Positional Encoding（原论文）}
\[
  \mathrm{PE}_{\mathit{pos},\, 2i}   \;=\; \sin\!\left(\frac{\mathit{pos}}{10000^{2i / d_{\text{model}}}}\right),
\]
\[
  \mathrm{PE}_{\mathit{pos},\, 2i+1} \;=\; \cos\!\left(\frac{\mathit{pos}}{10000^{2i / d_{\text{model}}}}\right).
\]
偶数维放 sin，奇数维放 cos，频率沿维度按几何级数变化：低维高频、高维低频。
\end{mathbox}

\begin{remark}[要不要学位置？]
也可以把 $\mathrm{PE}$ 改成\textbf{可学习的位置 embedding}（GPT 系就是这么做的）。
正弦版本的卖点是：它\textbf{不依赖训练长度}，理论上对超过训练时见过的序列长度也能给出有定义的位置向量；
学习版本表达力更强，但只在训练长度范围内有意义。两者在常规长度上效果差不多。
\end{remark}

\subsubsection{把整体拼起来}

万事俱备。一个完整的 Transformer 长这样：

\textbf{Encoder}（共 $N$ 层，原论文 $N=6$），每层：
\begin{enumerate}
  \item Self-Attention $\to$ Add \& Norm
  \item FFN $\to$ Add \& Norm
\end{enumerate}

\textbf{Decoder}（共 $N$ 层），每层：
\begin{enumerate}
  \item \textbf{Masked} Self-Attention $\to$ Add \& Norm \quad(\emph{只看自己前面已生成的位置})
  \item \textbf{Cross-Attention} $\to$ Add \& Norm \quad(\emph{Q 来自 decoder，K/V 来自 encoder 顶层输出})
  \item FFN $\to$ Add \& Norm
\end{enumerate}

\textbf{输出头}：decoder 最后一层 $\to$ Linear $\to$ $\softmax$ over vocab，逐位置预测下一个 token。

\begin{keypointblock}
\textbf{Cross-attention 是什么？}
它就是第二章 Bahdanau / Luong attention 在 Transformer 里的``\textbf{再实现}''：
还是一样的 $\softmax(QK^\top/\sqrt{d_k})V$，
只是\emph{Q 来自 decoder 当前层}（``我现在想要什么信息''），
而 \emph{K, V 来自 encoder 的最终输出}（``源句子里有哪些信息可拿''）。
所以原本 RNN 时代的``encoder--decoder attention'' 现在只是 Transformer 里\textbf{众多 attention 之一}，机制完全统一。
\end{keypointblock}

\subsection{这一章真正该记住的}
\label{subsec:ch03-takeaways}

\begin{enumerate}
  \item \textbf{纯 attention 就够了}：Transformer 不用 RNN、不用卷积，完全靠 attention + 线性层 + 归一化，
        证明了 ``recurrence'' 不是序列建模的必要条件。
  \item \textbf{并行训练}是它真正的杀手锏：一句话所有位置同时算，训练速度比 RNN 高得多。
  \item \textbf{Scaled dot-product attention}：核心公式 $\softmax(QK^\top/\sqrt{d_k})V$，$\sqrt{d_k}$ 是为了让 softmax 不饱和。
  \item \textbf{Masked self-attention} 让 decoder 在并行训练时仍然只能看过去，把``推理时逐词''和``训练时一次性''统一起来。
  \item \textbf{Multi-head} 不是变多变大，而是把同样大小的表示\textbf{切成几份}，并行学多种关注模式，再用 $W_O$ 融合。
  \item \textbf{四件脚手架}缺一不可：FFN 做位置内的非线性消化，残差连接保证深层可训，
        LayerNorm 稳住激活尺度，位置编码补回 attention 缺失的顺序信息。
  \item \textbf{Encoder = self-attn + FFN $\times N$；Decoder = masked self-attn + cross-attn + FFN $\times N$。}
        Encoder--decoder 之间的桥梁就是 cross-attention，本质和第二章那套 attention 是同一回事。
\end{enumerate}

\section{子词切分：字节对编码 BPE}

\subsection{章节定位与零基础该问的问题}

前三章我们一直假设“词表是给定的”：模型有一张固定的词典 $V$，输入输出都从里面选词。
但真实文本里永远会出现训练时没见过的拼写、人名、专业术语、罕见词。
本章讨论的就是\textbf{在固定词表的限制下，如何让模型仍然能处理“没见过的词”}。
答案是\textbf{子词（subword）切分}，最常见的实现就是 BPE。

\begin{problemblock}[零基础读者该追问的四个问题]
\begin{enumerate}
    \item 什么叫\textbf{固定词表问题（fixed vocabulary problem）}？模型为什么不能临时加词？
    \item 什么是 \textbf{UNK}？为什么它对翻译这种任务是灾难？
    \item BPE 的\textbf{训练阶段}到底在做什么？它学出来的是“词典”还是“规则”？
    \item BPE 的\textbf{推理阶段}又在做什么？面对一个全新的词，它怎么把它切开？
\end{enumerate}
\end{problemblock}

\subsection{固定词表问题}

\begin{problemblock}[问题]
神经网络的输入层 / 输出层维度在搭建时就已经写死：
输入 embedding 矩阵是 $|V|\times d$，输出 softmax 是 $d\times|V|$。
$|V|$ 一旦定下来，模型就只认识这 $|V|$ 个 token。
那遇到词表外的词怎么办？
\end{problemblock}

\begin{intuitionblock}[朴素直觉]
最朴素的办法有两端：
\begin{itemize}
    \item \textbf{词级（word-level）}：每个英文单词当一个 token。优点是“一个词一个意思”很自然；
          缺点是英文 + 形态变化 + 人名 + URL 几乎是无穷的，词表 50k 也兜不住，剩下的全部塞进一个特殊 token \texttt{<unk>}。
          一旦句子里出现 \texttt{<unk>}，翻译模型就只能输出 \texttt{<unk>}，相当于直接放弃。
    \item \textbf{字符级（char-level）}：把句子拆成一个个字符。词表只有几十到几千，永远不会 OOV；
          但代价是序列变得极长（一个 20 字母的英文词 $\to$ 20 个 token），
          注意力开销 $O(n^2)$ 爆炸，而且语义被打散到单个字母上，模型学得慢。
\end{itemize}
中间路线很自然：\textbf{常见词保持完整，罕见词切成更小的、训练时见过的子片段}。
这就是子词。
\end{intuitionblock}

\begin{remark}
BPE 被引入 NMT 的工作是 Sennrich, Haddow \& Birch (2016)，
\emph{Neural Machine Translation of Rare Words with Subword Units}。
原本 BPE 是一种数据压缩算法，他们把它“反过来用”：
不是为了压缩字节，而是为了\textbf{学一张折中粒度的词表}。
\end{remark}

\subsection{BPE 算法的两个阶段}

BPE 的核心只有一张东西——\textbf{有序合并表（merge table）}。
训练阶段\textbf{学这张表}，推理阶段\textbf{按这张表切词}。

\subsubsection{训练阶段：学合并规则}

\begin{mathbox}{BPE 训练循环}
\textbf{初始化}：
\begin{itemize}
    \item 词表 $V_0$ = 语料里出现过的\textbf{所有单字符}（外加一个词尾标记，例如 \texttt{</w>}）；
    \item 合并表 $M = \emptyset$；
    \item 把每个单词预先切成字符序列，例如 \texttt{cat</w>} $\to$ \texttt{c\,a\,t\,</w>}。
\end{itemize}
\textbf{重复 $N$ 次}（$N$ 是想要的合并次数，决定最终词表大小）：
\begin{enumerate}
    \item 在整个语料上数所有\textbf{相邻符号对}的频率 $\mathrm{count}(a, b)$；
    \item 取频率最高的那一对 $(a^\star, b^\star)$；
    \item 把规则 $a^\star + b^\star \to a^\star b^\star$ 追加到合并表 $M$；
          把新 token $a^\star b^\star$ 加进词表 $V$；
    \item 在所有已切分的词里，把出现的 $a^\star, b^\star$ 相邻处合并。
\end{enumerate}
\end{mathbox}

\begin{remark}[为什么需要词尾标记]
如果不加 \texttt{</w>}（或 SentencePiece 用的 \texttt{\_}）这种边界符号，
合并 \texttt{e+r} 会同时影响 \texttt{lower} 里的 \texttt{er} 和 \texttt{her}。
加了边界符号之后，\texttt{er</w>} 和 \texttt{er}（词中间的）就是两个不同符号，
保证一次合并只引入\textbf{一个}新 token。
\end{remark}

\begin{engbox}[实现上的小技巧]
真正实现时不会逐字符扫语料，而是：
\begin{itemize}
    \item 先把语料压成 \{单词 $\to$ 词频\} 的字典；
    \item 数“相邻符号对”时按词频加权，例如 \texttt{cat</w>} 出现 4 次，
          那么对 \texttt{(c,a)}、\texttt{(a,t)}、\texttt{(t,</w>)} 的贡献都各 +4；
    \item 这样每轮合并的复杂度只跟\textbf{不同单词的数量}有关，而不是语料的总长度。
\end{itemize}
\end{engbox}

\begin{example}[一个迷你语料]
设语料的单词频次表为
\[
\texttt{cat}\!\times\!4,\;
\texttt{mat}\!\times\!5,\;
\texttt{mats}\!\times\!2,\;
\texttt{mate}\!\times\!3,\;
\texttt{ate}\!\times\!3,\;
\texttt{eat}\!\times\!2.
\]
取 $N=5$ 次合并。初始词表是字符 $\{a,c,e,m,s,t,\texttt{</w>}\}$，每个词都切成字符。
执行训练循环时，最频繁的相邻对会按下面这种顺序被选出：
\begin{enumerate}
    \item \texttt{a+t} 出现在 \texttt{cat, mat, mats, mate, ate, eat}，是最频繁的对，
          合并 $\Rightarrow$ 新 token \texttt{at}；
    \item 接下来 \texttt{m+at}（来自 mat / mats / mate）频率领先，合并 $\Rightarrow$ \texttt{mat}；
    \item 再下来 \texttt{at+</w>}（来自原本的 cat / mat / ate）合并 $\Rightarrow$ \texttt{at</w>}；
    \item \texttt{c+at</w>} 合并 $\Rightarrow$ \texttt{cat</w>}；
    \item \texttt{e+at}（来自 eat）或者 \texttt{mat+e}（来自 mate）二选一被选出。
\end{enumerate}
五次合并后，常见整词如 \texttt{cat</w>}, \texttt{mat} 已经成为单 token，
而 \texttt{mate}、\texttt{mats} 之类则被表达成 2--3 个子词的组合。
精确顺序取决于实现的 tie-breaking 规则，但\textbf{大方向一定是“先合并最频繁的对”}。
\end{example}

\subsubsection{推理阶段：按合并表切新词}

\begin{mathbox}{BPE 推理（编码）}
对一个新词 $w$：
\begin{enumerate}
    \item 先把它切成字符序列（带词尾标记）；
    \item 在当前序列上，找出\textbf{所有可能的相邻合并}，
          其中“可能”指对应的规则在合并表 $M$ 中存在；
    \item 选\textbf{合并表里位置最靠前}的那条规则（位置越靠前 = 训练时越早被合并 = 越频繁），应用它；
    \item 重复 2--3 步，直到没有任何可应用的规则为止；
    \item 剩下的符号序列就是该词的 BPE 子词切分。
\end{enumerate}
\end{mathbox}

\begin{intuitionblock}[关键：合并表是\emph{有序}的]
训练时第 1 条规则学出来的 pair 一定比第 100 条更频繁。
推理时按这个顺序贪心地合并，相当于优先把“常见整块”凑出来。
所以同一个词在不同位置只会被切成同一种结果，BPE 的编码是\textbf{确定性}的。
\end{intuitionblock}

\begin{example}[已见词 vs.\ 未见词]
假设训练后合并表里靠前的几条是
$\texttt{a+t}, \texttt{m+at}, \texttt{at+</w>}, \texttt{c+at</w>}, \texttt{e+at}$。
\begin{itemize}
    \item 输入 \texttt{cat}：字符化为 \texttt{c\,a\,t\,</w>}
          $\to$ 合并 \texttt{at} $\to$ \texttt{c\,at\,</w>}
          $\to$ 合并 \texttt{at</w>} $\to$ \texttt{c\,at</w>}
          $\to$ 合并 \texttt{cat</w>} $\to$ 单 token \texttt{cat</w>}。\\
          \emph{常见词 = 一个 token。}
    \item 输入 \texttt{matador}（训练时没见过）：可能的子词切分是
          \texttt{mat\,a\,d\,o\,r\,</w>} —— \texttt{mat} 来自合并表，
          后面只能停在字符级。\emph{罕见词 = 一小串子词，但仍然没有 \texttt{<unk>}。}
\end{itemize}
\end{example}

\subsection{一些值得知道的工程细节}

\begin{engbox}[实现家族]
\begin{itemize}
    \item \textbf{原始 BPE}（subword-nmt）：依赖语言级别的预分词（按空格切词），
          在每个词内部跑 BPE，使用 \texttt{</w>} 标记词尾。
    \item \textbf{SentencePiece}：把空格也当成普通符号 \texttt{\_}（U+2581），
          \emph{不依赖任何语言相关的预分词}，因此天然支持中文/日文这种没有空格的语言。
          GPT-2/3、LLaMA、Qwen、Gemma 等大模型基本都用 SentencePiece 风格的 BPE。
    \item \textbf{YouTokenToMe / tiktoken / tokenizers (HF)}：对训练 / 编码做了大量工程优化，
          编码速度可以到数十 MB/s。
\end{itemize}
\end{engbox}

\begin{remark}[词表大小的经验值]
小任务的 BPE 词表通常 $4\text{k}\!\sim\!16\text{k}$；
NMT 经典设置 $32\text{k}$；
现代大模型 $50\text{k}\!\sim\!200\text{k}$（多语言往大走）。
词表越大，平均序列越短，但 embedding 矩阵越占显存。
\end{remark}

\begin{remark}[同家族的近亲]
\begin{itemize}
    \item \textbf{WordPiece}（BERT）：思路非常像 BPE，但选择合并对的准则不是“频率最高”，
          而是\emph{使语言模型似然增加最多}。
    \item \textbf{Unigram LM}（SentencePiece 默认选项之一）：
          反过来——先放一个大词表，再用 EM 慢慢删除“最没用”的子词。
    \item \textbf{BPE-dropout}：训练时以一定概率\emph{跳过某些合并}，
          于是同一个词会被随机切成多种合理子词序列，相当于一种数据增强，对低资源任务很有用。
\end{itemize}
\end{remark}

\subsection{这一章真正该记住的}

\begin{keypointblock}[Take-aways]
\begin{enumerate}
    \item 神经模型词表 $V$ 是固定的，词级切分必然遇到 OOV $\to$ \texttt{<unk>}，对生成任务尤其致命。
    \item 字符级永远没有 OOV，但序列太长、语义被打散；\textbf{子词是粒度上的折中}。
    \item BPE 的本质是一张\textbf{有序的合并表}：训练阶段按“相邻对频率最高”不断合并，逐条学出来。
    \item 推理阶段按这张表\textbf{从前往后贪心}地把字符序列合并起来；常见词收敛成一个 token，罕见词留成几段。
    \item 几乎所有现代 LLM 用的 tokenizer 都是 BPE 或它的近亲（WordPiece / Unigram LM），
          理解 BPE 就理解了它们 80\% 的行为。
\end{enumerate}
\end{keypointblock}

\section{分析与可解释性 (Analysis and Interpretability)}
\label{sec:analysis}

\subsection{章节定位 + 零基础该问的问题}

前面几章我们一路堆模型：先有 seq2seq，再加 attention，再加 multi-head，
每一步都伴随着一个“直觉理由”——比如“注意力应该能学到对齐”，
“多头应该能学到不同种类的关系”。
但这些都只是\emph{设计动机}，并不等于模型真的照做了。
这一章要做的事是：\textbf{打开训练好的模型，看看里面到底发生了什么}。

\begin{problemblock}{零基础读者会问的三个问题}
\begin{enumerate}
  \item 多头注意力的每个 head 真的在干不同的事吗？还是大家学的都差不多？
  \item 既然这样，能不能干脆把一些“没用”的 head 裁掉？
  \item Encoder 内部的向量表示，到底\emph{编码了什么语言学信息}（词性？形态？句法？）？
\end{enumerate}
\end{problemblock}

下文以两条研究线分别回应：第一条来自
Voita et al.\ (ACL 2019, \emph{Analyzing Multi-Head Self-Attention}),
第二条来自 Belinkov et al.\ (2017) 等围绕 NMT encoder 的 probing 工作。

\subsection{多头自注意力的头到底在干嘛？}
\label{sec:head-roles}

\subsubsection{重要的 head 是可解释的}

Voita et al.\ 训练好一个 Transformer NMT 模型后，逐个 head 分析其行为，
得到一个非常干净的结论：\textbf{真正“干活”的只是少数 head}，
而这些 head 的角色\emph{可以被人读懂}。

\begin{intuitionblock}
不是每个 head 都是平等的。一个训练好的多头自注意力，
更像是“几个专家 + 一堆陪跑”。
\end{intuitionblock}

被识别出的“专家”大体可以归成三类：

\begin{description}
  \item[(a) 位置型 head (positional heads)]
    几乎只看相邻位置。
    Encoder 中典型地有 2--3 个“看前一个 token”的 head，
    以及 \(\sim\)2 个“看后一个 token”的 head。
    这相当于模型自己学出了一个轻量级的局部卷积。
  \item[(b) 句法型 head (syntactic heads)]
    会跟踪“主语\(\to\)动词”“动词\(\to\)宾语”这类语法关系。
    这种行为在多个不同语言对的模型里都能观察到，
    说明它不是某种偶然，而是 multi-head 自然倾向于做的事。
  \item[(c) 稀有词型 head (rare-tokens head)]
    在第一层往往会出现一个特殊的 head：它把注意力集中到句子里最稀有的 token 上。
    一个合理的解释是，模型在“锚定”那些信息量最大的稀有线索；
    同时这个现象也带有一点轻微的过拟合气味
    （稀有 token 在训练集里本来就更具区分度）。
\end{description}

这个结果反过来验证了 multi-head 的归纳偏置：
\textbf{不同 head 确实会学到不同的东西}，并不是简单地把同一种关系算了 \(h\) 遍。

\subsubsection{大多数 head 可以被剪掉}

同一篇论文进一步做了一件激进的事：
给每个 head 套上一个可学习的“门”（gate），
通过稀疏正则把不重要的 head 推向 0，
然后看模型质量随被裁 head 数量的下降。

观察到两个互相印证的现象：

\begin{enumerate}
  \item \textbf{很大一部分 head 可以被裁掉}，翻译质量几乎不掉。
        典型设置下，只剩下 1/4 \(\sim\) 1/3 的 head，BLEU 仍然接近原模型。
  \item \textbf{留下来的 head 恰好就是上面归类的那些“专家”}
        （位置型、句法型、稀有词型）。
\end{enumerate}

但有一个看起来反直觉、却非常重要的细节：

\begin{keypointblock}
\textbf{这\emph{不}意味着可以一开始就训练一个 head 数很少的模型。}
直接用少头训练，质量明显更差。
多头在训练时的真正作用更像是\emph{一个更大的假设空间}：
让模型有机会“尝试”各种行为模式，最后挑出有用的那几个。
剪枝是事后压缩，而不是先验设计。
\end{keypointblock}

\subsection{Probing：表示里到底编码了什么？}
\label{sec:probing}

第二个问题指向 encoder 输出的那一堆向量：它们里面除了“能用来翻译”，
是否还藏着具体的语言学信息？衡量这件事的标准工具叫 \textbf{probing}（探针）。

\subsubsection{Probing 的方法论}

\begin{mathbox}{Probing 的标准流程}
\begin{enumerate}
  \item \textbf{冻结}已经训练好的 NMT encoder（参数不再更新）。
  \item 把一批带标注的数据（例如每个词都有词性标签）喂进去，
        在指定层取出每个词的向量表示 \(h_i\)。
  \item 在这些\emph{冻结}的 \(h_i\) 之上，训练一个\emph{小}分类器
        （通常是线性层或浅层 MLP）去预测目标标签 \(y_i\)。
  \item 用这个小分类器的准确率作为代理指标，
        近似回答“\(h_i\) 里有多少关于 \(y_i\) 的信息”。
\end{enumerate}
\end{mathbox}

\begin{remark}
一个常被忽视的注意点：probing 的准确率本身并不直接等于
“表示里编码了多少信息”，因为探针自己就有一定学习能力，
甚至能从随机特征里学出一些东西。
严格的 probing 研究通常会做对照实验（比如随机初始化的 encoder、
打乱标签的对照、或控制探针容量）。
本节只用其结论性观察，对方法论细节不再展开。
\end{remark}

\subsubsection{NMT encoder 学到了什么形态学？}

Belinkov et al.\ (2017) 用 probing 方法系统研究了 NMT encoder。
两个值得记住的结论：

\paragraph{(a) 词性 (POS) 与层数的关系}
拿 encoder 各层的表示去 probe POS 标注：
\begin{itemize}
  \item 训练后的 encoder 各层都比原始词嵌入显著更好；
  \item 但有趣的是，\textbf{第 1 层的 POS 探针准确率往往\emph{高于}第 2 层}。
\end{itemize}
对应的解释假设是：
第 1 层更接近“词级结构”（词性、形态等表层属性），
第 2 层开始向更抽象、更语义化的方向移动，
反而把表层语法信息“压”了下去。
这与“层数越深 = 越好”的朴素直觉是矛盾的。

\paragraph{(b) 目标语言的反向影响}
固定 source 语言（比如阿拉伯语），变化 target 语言，再 probe source 端形态信息：

\begin{intuitionblock}
直觉上，target 越“强”，encoder 越要努力工作，所以 encoder 的表示应该越丰富。
\end{intuitionblock}

实际观察到的是这样一个稍微反直觉的结果：

\begin{quote}
\textbf{Target 语言形态\emph{越贫乏}（如英语），encoder 对 source 形态的编码反而\emph{越好}；
target 语言形态\emph{越丰富}（如希伯来语、德语），encoder 反而编码得更弱。}
\end{quote}

一个解释方式是：
当 target 自己有丰富的屈折变化时，模型可以“偷懒”，
让大部分形态信息直接从 target 端的解码器/词表那里产生；
而 target 一旦贫乏（英语几乎没有屈折），
encoder 就不得不在自己这一侧把 source 的形态彻底搞清楚，
否则解码器没法做出正确翻译。

\subsection{这一章真正该记住的}

\begin{keypointblock}
\begin{enumerate}
  \item \textbf{头不平等}：训练好的多头注意力里，只有一小部分 head 真正承担主要工作。
  \item \textbf{角色可解释}：重要的 head 倾向于扮演位置型、句法型、稀有词型这几种角色，
        说明 multi-head 的归纳偏置确实在起作用。
  \item \textbf{剪枝可行，但前提是先有多头}：可以事后裁掉大部分 head 而几乎不掉点；
        但直接用少头从头训练效果会明显更差——多头主要是\emph{训练阶段的假设空间}。
  \item \textbf{Probing 是“表示里有什么”的标准工具}：
        冻结 encoder，在表示之上训练浅层分类器，用准确率近似信息量；
        但要小心探针自身能力带来的偏差。
  \item \textbf{Target 越弱，encoder 越强}：
        当 target 语言形态贫乏时，NMT encoder 反而会把 source 的形态信息编码得更彻底。
\end{enumerate}
\end{keypointblock}


\section*{致谢与原文}
\addcontentsline{toc}{section}{致谢与原文}
本文所有图示思路、公式与例子来源于
Lena Voita 的公开课程页面：
\url{https://lena-voita.github.io/nlp_course/seq2seq_and_attention.html}。
本中文讲义只是按零基础读者的视角做了重写与展开，
以便把原页面的「先看图、再看公式」体验补回到纯文本场景。

\end{document}
